\section{Justice and Mercy United in Christ the Judge}

A shallow imagination splits God into two rival faces: a stern Judge enforcing justice and a gentle Savior pleading mercy. Scripture knows no such division. The Psalmist joins both in one song: ``Mercy and judgment I will sing to thee, O Lord'' (Ps 100:1)---the psalm Evodius raised at Ostia when St.\ Monica died.\endnote{St.\ Augustine, \textit{Confessions} IX.12.31, in \textit{NPNF} I, vol.\ 1.} Psalm 84 beholds them embracing: ``Mercy and truth have met each other: justice and peace have kissed'' (Ps 84:11). Psalm 144 proclaims that ``his tender mercies are over all his works'' (Ps 144:8--9), and St.\ James binds warning and promise together: ``For judgment without mercy to him that hath not done mercy. And mercy exalteth itself above judgment'' (James 2:13).

St.\ Thomas Aquinas establishes the theological root of this harmony in Question 21 of the \textit{Prima Pars}.\endnote{St.\ Thomas Aquinas, \textit{Summa Theologiae} I, q.\ 21, aa.\ 1--4.} Divine justice is not commercial exchange, for ``who hath first given to him, and recompense shall be made him?'' (Rom 11:35); it is distributive order conforming every creature to the eternal rule of God's wisdom and goodness. Hence ``God's justice, which establishes things in the order conformable to the rule of His wisdom, which is the law of His justice, is suitably called truth.''\endnote{Aquinas, \textit{ST} I, q.\ 21, aa.\ 1--2.}

To the objection that mercy would contradict justice by relaxing God's sentence, St.\ Thomas replies:

\begin{quote}
God acts mercifully, not indeed by going against His justice, but by doing something more than justice; thus a man who pays another two hundred pieces of money, though owing him only one hundred, does nothing against justice, but acts liberally or mercifully. The case is the same with one who pardons an offence committed against him, for in remitting it he may be said to bestow a gift. Hence the Apostle calls remission a forgiving: ``Forgive one another, as Christ has forgiven you'' (Eph.\ 4:32). Hence it is clear that mercy does not destroy justice, but in a sense is the fulness thereof. And thus it is said: ``Mercy exalteth itself above judgement'' (James 2:13).\endnote{Aquinas, \textit{ST} I, q.\ 21, a.\ 3, ad 2.}
\end{quote}

Mercy never overthrows justice; it fulfills and surpasses it. In Article 4, St.\ Thomas proves further that mercy and justice unite in \textit{every} work of God, with mercy at the primary root:

\begin{quote}
Now the work of divine justice always presupposes the work of mercy; and is founded thereupon. For nothing is due to creatures, except for something pre-existing in them, or foreknown. Again, if this is due to a creature, it must be due on account of something that precedes. And since we cannot go on to infinity, we must come to something that depends only on the goodness of the divine will---which is the ultimate end. \dots\ So in every work of God, viewed at its primary source, there appears mercy. In all that follows, the power of mercy remains, and works indeed with even greater force; as the influence of the first cause is more intense than that of second causes.\endnote{Aquinas, \textit{ST} I, q.\ 21, a.\ 4, corpus.}
\end{quote}

Every reward rendered in justice rests upon the unmerited mercy of creation, Incarnation, and grace (Luke 7:47); even in condemning the impenitent, mercy endures by punishing less than strict desert could claim.\endnote{Aquinas, \textit{ST} I, q.\ 21, a.\ 4, ad 1 and ad 3.} Purgatory itself mirrors this union: justice insists that nothing defiled enter heaven and every disordered debt be mended, while mercy refuses to discard a weak friend of God and provides both cleansing fire and the Church's suffrages.

Where did justice and mercy meet visibly in history? On Calvary. Commenting on Psalm 84:10--11, St.\ Robert Bellarmine notes that before Christ, mercy and justice seemed opposed---``the one prompting to punish, the other to forgive; for Christ's passion and suffering was meant to deliver the human race in mercy, while it made the fullest satisfaction to the divine justice''---so that on the Cross they joined ``in the bonds of the strictest friendship.''\endnote{St.\ Robert Bellarmine, \textit{Commentary on the Book of Psalms}, on Ps 84:10, trans.\ O'Sullivan (1866).} Yet from the next verse---``Truth is sprung out of the earth: and justice hath looked down from heaven'' (Ps 84:12)---Bellarmine draws a sobering warning:

\begin{quote}
For the extent of God's anger and hatred of sin would never have been thoroughly known, had not God decreed that it should be expiated by the death of his only Son; and, even, we should never have known the extent of God's anger to the sinner on the day of judgment, had we not seen the amount and the extent of Christ's sufferings in atoning for the sins of others, ``For if in the green wood they do these things, what shall be done in the dry?'' says our Lord, Luke 23.\endnote{St.\ Robert Bellarmine, \textit{Commentary on the Book of Psalms}, on Ps 84:11 (Ps 84:12 in Challoner numbering), trans.\ O'Sullivan (1866).}
\end{quote}

Calvary forbids presumption and despair together. Looking at the Cross, no one can treat sin lightly when it cost the blood of God's Son; looking at the same Cross, no penitent can doubt the Judge's mercy, since the Judge is the Redeemer who bled for us. St.\ John anchors our hope in that identity: ``if any man sin, we have an advocate with the Father, Jesus Christ the just. And he is the propitiation for our sins'' (1 John 2:1--2). St.\ Paul asks triumphantly: ``Who is he that shall condemn? Christ Jesus that died: yea that is risen also again, who is at the right hand of God, who also maketh intercession for us'' (Rom 8:33--34). And St.\ Alphonsus Liguori places St.\ Thomas of Villanova's prayer on our lips: ``I willingly submit to the judgment of him who died for me, and who, that he might not condemn me, has condemned himself to the cross.''\endnote{St.\ Alphonsus Maria de Liguori, \textit{Preparation for Death}, Consideration XXV, First Point, ``Affections and Prayers,'' trans.\ Duffy (1843), p.\ 304.}

Held in this balance, Catholic teaching guards the soul against two mortal distortions: presumption and despair.

\textit{Presumption} exploits Christ's mercy while ignoring His holiness, postponing repentance and clinging to occasions of sin. St.\ Gregory the Great shatters that complacency in Book XIV of his \textit{Moralia in Job}:

\begin{quote}
For the terribleness of the Judge cannot be avoided saving before the Judgment. Now He is not discerned, but is appeased by prayers. But when He shall sit on that dreadful inquest, He is both able to be seen and not able any longer to be propitiated; in that the doings of the wicked which He bore long while in silence, He shall pay back all of them together in wrath. Whence it is necessary to fear the Judge now, while He does not yet execute judgment, while He bears patiently for long, while He still tolerates the wickedness that He sees.\endnote{St.\ Gregory the Great, \textit{Moralia in Job} XIV.58.79, Library of the Fathers translation.}
\end{quote}

St.\ Augustine warns likewise in \textit{The City of God} that ``in vain, after this life, does a man seek for what he has neglected to provide while in this life.''\endnote{St.\ Augustine, \textit{The City of God} XXI.24, trans.\ Dods (1871).} And St.\ James reminds us that whoever refuses mercy to his neighbor saws off the branch of his own hope (James 2:13).

\textit{Despair} makes the opposite error, staring at divine strictness while looking away from Christ's wounds. It hears Psalm 129:3 (``If thou, O Lord, wilt mark iniquities: Lord, who shall stand it'') and Psalm 142:2 (``enter not into judgment with thy servant''), yet forgets the next breath: ``For with thee there is merciful forgiveness \dots\ Because with the Lord there is mercy: and with him plentiful redemption'' (Ps 129:4, 7--8). Against every despairing whisper, St.\ Thomas Aquinas teaches firmly: ``to say that in this life there is any sin of which one cannot repent, is erroneous, first, because this would destroy free-will, secondly, because this would be derogatory to the power of grace \dots\ It is also erroneous to say that any sin cannot be pardoned through true Penance.''\endnote{Aquinas, \textit{ST} III, q.\ 86, a.\ 1, corpus.} While breath lasts, the confessional remains open and Christ's blood outstrips every sin.

Between presumption and despair runs the narrow road of filial vigilance: ``But if we would judge ourselves, we should not be judged. But whilst we are judged, we are chastised by the Lord, that we be not condemned with this world'' (1 Cor 11:31--32). Judging ourselves now takes three daily forms: honest self-examination and regular Sacramental Confession, where self-accusation brings pardon rather than sentence; concrete works of mercy toward the poor and the Holy Souls (Matt 25:34--40; James 2:13); and patient endurance of earthly crosses united to Christ's Passion, letting grace burn away our wood, hay, and stubble before the grave.

Our Lord sums up this watchfulness with steadying hope: ``Watch ye, therefore, praying at all times, that you may be accounted worthy to escape all these things that are to come and to stand before the Son of man'' (Luke 21:34, 36). To stand before the Son of man unafraid belongs not to the proud, but to the penitent who has already met the Judge beneath His Cross.
